\section{A pressure-free sum of physical critical sectors}
\label{sec:mc-critical}
We now use the endpoint kernel on its own physical source, rather than
multiply a local critical estimate by a count of collision words. The
conclusion is a complete, collision-uniform leading-coefficient sum.
It removes a genuine pressure dependency of
Proposition~\ref{prop:cr-critical-sum}; it is not yet a normalized
roof-density theorem. All endpoint coordinates below are the original
angles $q=(\alpha_0,\alpha_m)$.

\begin{lemma}[Endpoint Hessian bounds and the momentum inverse weight]
\label{lem:mc-hessian}
Let $x$ be a critical point of a regular reduced action $F$ with every
visited incidence positive. Write
\[
 G=D_q^2F(x)=\begin{pmatrix}a_0&e\\e&d_0\end{pmatrix},
 \qquad \Delta=\det G>0,\qquad
 u_x=\frac1{|\partial_{\alpha_0}p_m(x)|}.
\]
Then $p_0=p_m=0$, and
\begin{align}
 u_x&=\frac{R|e|}{\Delta}=\frac{R^2\omega(x)}{\Delta},
   &\mathfrak J(x)&=\frac{\sqrt\Delta}{2R^2c_*}\,u_x,
                                                    \label{eq:mc-weight}\\
 \frac{\sqrt\Delta}{R^2}&\le\kappa+t_*^{-1}\le U=170/9,
   &\frac{\lambda_{\max}(G)}{R^2}
       &\le\kappa+2t_*^{-1}\le U_2:=320/9.\label{eq:mc-Hbounds}
\end{align}
For $m\ge2$ the second bound improves to
$\lambda_{\max}(G)/R^2\le\kappa+t_*^{-1}$.
The same auxiliary-action identities hold at a one-sided unvisited
boundary when all \emph{visited} incidences stay positive. No physical
$DT^m$ at a grazing collision is thereby defined.
\end{lemma}
\begin{proof}
The endpoint first variations in the retained optical proof are
$F_{\alpha_0}=-Rp_0$, $F_{\alpha_m}=Rp_m$. The generating-function
matrix in Proposition~\ref{thm:cr-monodromy} therefore has lower-left
entry $-\Delta/(Re)$. Its inverse absolute value gives the first
identity, since $\omega=|e|/R$; the retained critical coarea coefficient
$\mathfrak J=\omega/(2c_*\sqrt\Delta)$ gives the second.

Here are the upper bounds, including their signs and units. For one
physical flight of length $\tau$, differentiate its Euclidean length
with respect to the first contact angle. The diagonal second derivative
is $R c_0+R^2c_0^2/\tau$: the circular acceleration is $-Rn_0$
and its scalar product with the initial length gradient $-v$ is
$Rc_0$. At the terminal contact the incoming velocity has normal
component $-c_m$, giving the same positive circular term $Rc_m$.
At a critical orbit $c_0=c_m=1$. Consequently the two endpoint
diagonals of the unreduced full length action are at most
$R+R^2/t_*$. Eliminating the positive interior optical Hessian is a
Schur complement, which subtracts a positive semidefinite matrix from
the endpoint block. Thus both reduced diagonals are at most this value,
and $\sqrt\Delta\le\sqrt{a_0d_0}\le R+R^2/t_*$.
If $m\ge2$, the unreduced endpoint block is diagonal, so its maximum
eigenvalue has the same upper bound. If $m=1$ there is no interior to
eliminate: the mixed second derivative of the one-flight length has
absolute value at most $R^2/\tau$, and the two-row bound is
$R+2R^2/t_*$. Division by $R^2$ proves \eqref{eq:mc-Hbounds}.
The positivity and elimination are exactly the physical optical
Hessian construction retained in Section~\ref{sec:v109-action};
no full-trajectory Hessian is replaced by an independent matrix model.
At a skipped-unvisited tangency its strictly reflected auxiliary word
is still a smooth action with positive visited incidences. The same
calculation is valid for that auxiliary action, restricted to its
actual physical side as in Section~\ref{sec:cr-critical}.
\end{proof}

\begin{theorem}[Complete regular critical coefficient sum]
\label{thm:mc-critical-regular}
For every $R$, every $m\ge1$ and every exact record $(n,k)$, the
one-copy sum over all regular critical orbits of the original $m$-flight
roof in that record satisfies
\begin{equation}\label{eq:mc-critical-regular}
 \sum_x\mathfrak J(x)\le\frac{U C_{\rm mom}}{2c_*}.
\end{equation}
This also bounds the sum with any coefficients between zero and one,
including restrictions to original roof intervals. Neither the
whole-cylinder pressure input nor an independent periodic-orbit
counting theorem is used. Critical orbits are the actual normal-endpoint
reversing trajectories identified in the retained proof, not periodic
points of a replacement symbolic model.
\end{theorem}
\begin{proof}
At every such $x$, $p_0=p_m=0$. Distinct physical initial points are
distinct regular roots in \eqref{eq:mc-roots} at $(a,v)=(0,0)$.
The pointwise root bound, rather than an arbitrarily chosen version
of an almost-everywhere density, gives $\sum_x u_x\le C_{\rm mom}$.
Use \eqref{eq:mc-weight}--\eqref{eq:mc-Hbounds} and sum nonnegative
terms. Exact section and record restrictions can only decrease the
sum. The proof works for the complete countable collection, by finite
exhaustion, before taking the supremum over radius or collision count.
\end{proof}

\subsection{One-sided sectors are paid on the physical side}
Bulk singular-nullness alone proves nothing about an endpoint trace.
We give the additional argument for the one-sided critical strata used
in the retained local coarea theorem. Consider a finite or countable
one-copy family of critical auxiliary-action germs $x$. Every germ has
all visited incidences positive and normal endpoints. Its physical-side
set $E_x$ retains all unvisited-obstacle inequalities and the original
section and record cuts. In a neighborhood of $x$, assume these cuts
are a finite Boolean combination of $C^1$ inequalities whose active
boundary differentials are nonzero. The limiting tangent set is then
a finite union of angular sectors. Multiple simultaneous boundaries
are permitted; no positive trace is discarded.

This class includes the original finite-time section/record cuts on a
regular auxiliary word. A section boundary is a coordinate boundary
pulled back by a locally invertible finite collision map, so its
boundary differential is nonzero. The exact occupation count is a
finite Boolean combination of these tests. An unvisited tangency is
a clearance boundary: the line-coordinate determinant in
Lemma~\ref{thm:cf-flux} makes its differential nonzero, and
an interior-foot tangency stays away from the incident contacts.
Only finitely many possible centers occur for a finite word by finite
horizon. Strict displacement labels are locally constant on such a
word. General Borel pairing selectors, or strata involving an actual
visited grazing collision, are \emph{not} asserted to have this tangent
sector property; their exact local selectors and outstanding budgets
remain in the manuscript.

Put $y=G_x^{1/2}(q-q_x)$ and let $\beta_y(x)$ be the angular fraction
of the physical tangent set in these normalized coordinates. It is
the leading coefficient fraction in the local theorem of
Section~\ref{sec:cr-critical}. Let $z=(p_0,p_m)$ and denote the
corresponding angular fraction in the $z$-plane by $\beta_z(x)$.
Auxiliary branches are identified by their strict visited word; their
source restrictions are disjointized on overlaps. No multiple copy is
introduced at a coordinate seam. The family below may also contain
fully regular critical points, whose interior sector fraction is one.

\begin{lemma}[The correct angular change of variables]
\label{lem:mc-angle}
Let $D=\operatorname{diag}(-1,1)$. At a critical point,
\begin{equation}\label{eq:mc-linear}
 z=R^{-1}D G_x^{1/2}y+o(|y|),\qquad L_x=R^{-1}D G_x^{1/2}.
\end{equation}
Writing $e_\theta=(\cos\theta,\sin\theta)$, the absolute angular
Jacobian is
\begin{equation}\label{eq:mc-angle}
 \left|\frac{d\arg(L_xe_\theta)}{d\theta}\right|
       =\frac{\sqrt{\det G_x}}{e_\theta^{\mathsf T}G_xe_\theta}.
\end{equation}
For every such tangent sector set,
\begin{equation}\label{eq:mc-sector-comparison}
 \mathfrak J(x)\beta_y(x)
       \le\frac{\lambda_{\max}(G_x)}{2R^2c_*}
                                     u_x\beta_z(x).
\end{equation}
The Euclidean angle in the unnormalized endpoint coordinates is not
substituted for either fraction.
\end{lemma}
\begin{proof}
Use $g=\nabla_qF=G_x(q-q_x)+o(|q-q_x|)$ and $z=R^{-1}Dg$.
For an invertible planar matrix $L$, differentiation of its polar angle
gives $|\det L|/|Le_\theta|^2$. Here
$|\det L_x|=\sqrt{\det G_x}/R^2$ and
$|L_xe_\theta|^2=e_\theta^{\mathsf T}G_xe_\theta/R^2$.
Integrating \eqref{eq:mc-angle} over the sector yields
$\beta_z\ge\sqrt\Delta\,\beta_y/\lambda_{\max}(G_x)$.
Combine this with \eqref{eq:mc-weight}. This argument also permits
finite unions, half-planes, corners and complements; boundary rays
have zero angular measure.
\end{proof}

\begin{theorem}[Complete one-sided leading critical budget]
\label{thm:mc-critical-sectors}
For the physical-side family just specified, with unique strict
visited words and one-copy source restrictions, one has
\begin{equation}\label{eq:mc-sector-momentum}
 \sum_x u_x\beta_z(x)\le C_{\rm mom},
\end{equation}
and consequently, uniformly in $R$, $m\ge1$, exact labels and any
restriction of the critical values to an original roof interval,
\begin{equation}\label{eq:mc-critical-sectors}
 \sum_x\mathfrak J(x)\beta_y(x)
                       \le\frac{U_2 C_{\rm mom}}{2c_*}.
\end{equation}
Here $u_x$ at an unvisited boundary is the auxiliary inverse weight
$R|e|/\Delta$, not a claimed derivative of the singular physical map.
This conclusion removes the whole-cylinder pressure premise for this
complete leading critical-sector sum. It does not extend a physical
trajectory through a grazing reflection or assert that every Borel
pairing seam possesses a sector limit.
\end{theorem}
\begin{proof}
First take a finite collection $\mathcal F$ of distinct physical auxiliary germs.
A strictly reflected finite word with positive visited incidences has
an invertible endpoint-momentum map: its endpoint derivative has
nonzero determinant, as follows from $e\ne0$ and $\Delta>0$.
The inverse-function theorem gives a common small $z$-disk for the
finitely many local inverses. The absolute Jacobian from
$(p_0,p_m)$ to $(p_0,\alpha_0)$ on germ $x$ is a positive continuous
function $u_x(z)$ with $u_x(0)=u_x$.

Distinct regular initial points have disjoint sufficiently small
source neighborhoods. The same is true for distinct strict-word
unvisited-boundary germs: at a limiting state, continuing straight
past its unvisited tangencies fixes the next strictly met disk
uniquely; iterate this observation through the fixed $m$ strict
collisions. Any alternative word using one of the tangencies would
have a visited zero incidence and is outside this class. Thus two
such germs at the same initial state are the same germ; their
selectors are united, not counted twice. This is a uniqueness
argument for the auxiliary extension and does not alter a physical
collision record. Artificial partition seams are identified as well.
Finite disjointization of actual source restrictions handles any
remaining coordinate overlap.

On the actual side of each germ, its auxiliary orbit is the actual
regular physical orbit. Therefore its positive contribution to the
joint endpoint root density is bounded, almost everywhere in $z$, by
the complete kernel already proved:
\begin{equation}\label{eq:mc-local-domination}
  \sum_{x\in\mathcal F}u_x(z)\one_{E_x}(z)
                      \le C_{\rm mom}.
\end{equation}
No equality on the singular state itself is needed. By polar Fubini,
the inequality holds at almost every angle for almost every radius
$r>0$. Choose a sequence of such radii tending to zero. On each
fixed finite family, the $C^1$ inequalities with nonzero active
linear parts imply convergence of $\one_{E_x}(re_\theta)$ to its
momentum tangent-sector indicator for almost every $\theta$.
Continuity of $u_x(z)$ and bounded convergence on the angular circle
give, after dividing the integral of \eqref{eq:mc-local-domination}
by $2\pi$,
$\sum_{x\in\mathcal F}u_x\beta_z(x)\le C_{\rm mom}$.
Take the supremum over finite subfamilies to prove
\eqref{eq:mc-sector-momentum}. Countability causes no exchange of a
signed series or of the collision and outer limits. Finally apply
Lemma~\ref{lem:mc-angle} and \eqref{eq:mc-Hbounds}. Keeping original
record, critical-roof or bounded continuous guard weights only reduces
the nonnegative sum. Neither independence nor a physical-word
multiplicity factor is used.
\end{proof}

\subsection{Exact gain and the remaining normalized obstruction}
Theorems~\ref{thm:mc-critical-regular} and
\ref{thm:mc-critical-sectors} prove a complete leading-coefficient
bound at every collision count, including genuine one-sided physical
critical sectors. Their source is the full original momentum
coordinate disintegration, and their constants are uniform on the
original radius family. This replaces the pressure-dependent
coefficient estimate, not its historical proof body.

There is no illicit interchange of a local-radius limit with $m\to\infty$.
The sector proof first treats finitely many germs at a \emph{fixed}
$R,m$, takes the local radius to zero, and exhausts their positive
coefficients. The resulting bound is uniform in $R,m$ because its
right-hand side is. It does not give a common neighborhood size or a
summable bound for the local $O(\sqrt s)$ errors. In particular the
critical values of different germs can differ, so the bound at the
common momentum origin cannot be used as a bound at different
momentum points selected by a common physical roof level.

The full original problem still requires the unchanged limit
\eqref{eq:uc-open-limit}, with $B$-dependent widths fixed before the
collision limsup. Along the present route the minimal additional
inputs are: a complete original-roof chart/error estimate with the
required $m^{-2}$ normalization, including the first-order clearance
source, noncritical folds and actual visited singular strata; payment
of the original incidence-only low-load source or a valid full
nongrazing-capacity comparison; and the stated spectral, arithmetic,
pairing and conditional-denominator inputs with their actual uniform
quantifiers. General pairing selectors retain their exact local
angular functions unless their tangent sectors are proved to exist.
Neither the nonzero single-mark mass nor a positive boundary trace
is set to zero. The complete raw pointwise local limit and final
four-coordinate inversion are not asserted by the present theorems.
